\begin{trapbox}

\begin{description}[style=nextline, leftmargin=2.8em, labelsep=0.5em, itemsep=0.35em]

    \item[\textbf{\textcolor{CTrap}{易错点一（混淆独立与互斥）}}]
    误认为两个事件互斥，就可以直接按照独立事件使用
    \[
        P(AB)=P(A)P(B).
    \]

    \item[\textbf{\textcolor{CTrap}{正确理解（区分定义）：}}]
    若 $A,B$ 互斥，则
    \[
        P(AB)=0.
    \]
    若 $A,B$ 又相互独立，则还必须满足
    \[
        P(A)P(B)=0,
    \]
    即至少有一个事件的概率为 $0$。

    因此，当
    \[
        P(A)>0,\qquad P(B)>0
    \]
    时，互斥事件一定不独立。

    \item[\textbf{\textcolor{CTrap}{错因分析（概念混淆）：}}]
    “互斥”表示两个事件不能同时发生；“独立”表示一个事件是否发生不影响另一个事件发生的概率。二者含义不同，不能混为一谈。

\end{description}

\medskip

\begin{description}[style=nextline, leftmargin=2.8em, labelsep=0.5em, itemsep=0.35em]

    \item[\textbf{\textcolor{CTrap}{易错点二（直接套用公式）}}]
    遇到 $P(AB)$ 就直接写成 $P(A)P(B)$，没有检查事件是否独立。

    \item[\textbf{\textcolor{CTrap}{正确理解（先判独立）：}}]
    使用
    \[
        P(AB)=P(A)P(B)
    \]
    计算积事件概率之前，必须确认 $A,B$ 相互独立，或者题目已经给出足以推出独立性的条件。

    \item[\textbf{\textcolor{CTrap}{错因分析（条件忽视）：}}]
    乘积公式并非对任意两个事件都成立；忽略独立性条件，会导致计算错误。

\end{description}

\medskip

\begin{description}[style=nextline, leftmargin=2.8em, labelsep=0.5em, itemsep=0.35em]

    \item[\textbf{\textcolor{CTrap}{易错点三（多个事件独立误区）}}]
    认为三个事件两两独立，就一定可以直接使用三个事件的乘积公式。

    \item[\textbf{\textcolor{CTrap}{正确理解（两两独立不够）：}}]
    对三个事件 $A_1,A_2,A_3$，相互独立不仅要求
    \[
        P(A_1A_2)=P(A_1)P(A_2),
    \]
    \[
        P(A_1A_3)=P(A_1)P(A_3),
    \]
    \[
        P(A_2A_3)=P(A_2)P(A_3),
    \]
    还必须满足
    \[
        P(A_1A_2A_3)
        =
        P(A_1)P(A_2)P(A_3).
    \]

    \item[\textbf{\textcolor{CTrap}{错因分析（条件遗漏）：}}]
    两两独立只检查了任意两个事件之间的关系，没有保证三个事件同时发生时仍满足乘积规律。

\end{description}

\end{trapbox}